\documentclass[10pt,a4paper]{amsart}
\usepackage[margin=19mm]{geometry}
\usepackage{amsmath,amssymb,mathtools}
\usepackage[scr=rsfs,cal=euler]{mathalfa}
\usepackage{xfrac}
\usepackage[unicode,hidelinks,bookmarks=false]{hyperref}
\newcommand{\ie}{i.e.\ }
\newcommand{\cf}{cf.\ }
\newcommand{\resp}{respectively\ }
\newcommand{\Subsection}{\subsection}
\providecommand{\nobreakdash}{\nobreak\hbox{-}}
\setlength{\emergencystretch}{2em}
\multlinegap0pt
\usepackage{mathtools}
\usepackage{amssymb}
\newtheorem{assumption}{Assumption}
\newtheorem{lemma}{Lemma}
\def\cal{\mathcal}
\title{Mean field social optimization: feedback person-by-person optimality and the dynamic programming equation}
\author{Minyi Huang, Shuenn-Jyi Sheu,, Li-Hsien Sun}
\date{Source: arXiv:2508.14236v2 (CC BY 4.0)}
\begin{document}
\maketitle

\noindent\emph{Source and licence.} Mean field social optimization: feedback person-by-person optimality and the dynamic programming equation. Version arXiv:2508.14236v2 (CC BY 4.0), distributed by arXiv under the Creative Commons Attribution 4.0 International licence, http://creativecommons.org/licenses/by/4.0/, as recorded in the arXiv licence metadata for this exact version and shown on the abstract page https://arxiv.org/abs/2508.14236v2. Full licence notice: \texttt{licenses/arxiv-2508.14236-CC-BY-4.0.txt}.\par\smallskip

\noindent\textit{Editorial scope.} Counted unit: the article's lemma lemma:xi (source0.tex lines 1093--1103). The article's complete original proof of it is delivered with it (source0.tex lines 1105--1168, one \texttt{proof} environment of 64 lines). No mathematics is written, repaired or re-derived: the statement and the proof are the article's own lines, in the article's own notation, and no retained line is substituted or re-flowed.  Retained uncounted as the source-local context the proof needs: lines 343--358; lines 445--447; lines 472--478; lines 524--536; lines 662--665; lines 2393--2396.  Nothing was omitted from any retained span, so the package carries no omission record.  No retained line contains a citation, so the wrapper carries no bibliography and no bibliography entry.  No retained line contains a figure environment, an image-inclusion command or drawing code, so no image is delivered and the package declares no asset dependency.  Editorial formatting only, in addition: an AMS \texttt{amsart} preamble that restates the article's own declarations and macros used by the retained lines, namely the environments \texttt{assumption}, \texttt{lemma} on separate counters, together with the standard packages those lines need.  The retained environments print in this document as Assumption 1, Lemma 1 in order of appearance, whereas the article prints the same items with the numbering of its own full document.  The complete native source of the article as distributed in the arXiv e-print for this exact version is bundled unchanged as \texttt{sources/183535/source0.tex}.
\begin{align}
- \partial_t V(t, x, \mu)
= \ &   V_x (t, x, \mu)
f(x, \hat u^i, \mu) + \frac{1}{2} {\rm Tr}
[ V_{xx}(t, x, \mu) (\Sigma+\Sigma_0)  ]
\label{MEVwoCN}   \\
&+ L(x, \hat u^i, \mu)  + \langle \mu(dy),  \partial_y\delta_\mu V(t,x, \mu; y) f( y, \phi(t, y, \mu),\mu)\rangle \notag  \\
&+\frac{1}{2} \langle \mu(dy), {\rm Tr}[
\partial_{y}^2 \delta_\mu
V(t, x, \mu; y) (\Sigma+\Sigma_0)]\rangle \notag \\
&+\langle \mu(dy),  {\rm Tr} [ \partial_{xy}\delta_\mu V(t,x, \mu;y)  \Sigma_0 ]\rangle \notag  \\
&+\frac{1}{2}\langle \mu^{\otimes 2}(dydz),  {\rm Tr}[\partial_{yz}\delta_{\mu\mu} V (t, x, \mu;y, z)
\Sigma_0]\rangle ,\notag\\
&\hskip -1cmV(T, x, \mu)=g(x,\mu),\qquad (t,x,\mu)\in[0,T]\times \mathbb{R}^{n}\times
 {\cal P}_2(\mathbb{R}^n), \label{vtcg}
\end{align}

\begin{assumption}\label{assum:meV}
The master equation \eqref{MEVwoCN} has a solution pair $(V, \phi)$ with $V\in {\cal C}_V$.
\end{assumption}

\begin{align}
dX_t^1=\ & f(X_t^1, u_t^1, \mu_{t}^{-1})dt+
\sigma dW_t^1+\sigma_0 dW_t^0,  \label{X1u} \\
dX_t^k=\ & f(X_t^k, \phi(t, X_t^k,\mu_t^{-k}), \mu_{t}^{-k})dt+
\sigma dW_t^k +\sigma_0 dW_t^0 ,   \label{Xkuhat} \\
& 0\le t\le T, \quad 2\le k\le N.   \notag
\end{align}

\begin{assumption}\label{assum:flg}
Each of the functions $ \delta_\mu f^*$, $ \delta_{\mu\mu} f^*$, $ \delta_\mu L^*$, $ \delta_{\mu\mu} L^*$, $ \delta_\mu g$ and $ \delta_{\mu\mu} g$ is jointly continuous in its arguments, and
there exists a constant $C_a$ such that
\begin{align*}
&|\delta_\mu f^*(t, x, \mu;y)| \le C_a(1+|x|+ |y|+\langle |y| \rangle_{\mu}  ),   \\
& |\delta_{\mu\mu} f^*(t, x, \mu;y,z)| \le C_a(1+|x|+ |y|+|z|+\langle |y| \rangle_{\mu}   ),    \\
 &|\delta_{\mu} L^*(t, x, \mu; y) |\le C_a(1+|x|^2+ |y|^2+\langle |y| \rangle_{\mu}^2),\\
&|\delta_{\mu\mu} L^*(t, x, \mu; y,z) |\le C_a(1+|x|^2+ |y|^2+|z|^2+\langle |y| \rangle_{\mu}^2),\\
&|\delta_{\mu} g( x, \mu; y) |\le C_a(1+ |x|^2+|y|^2+\langle |y| \rangle_{\mu}^2),\\
& |\delta_{\mu\mu} g( x, \mu; y,z) |\le C_a(1+ |x|^2+|y|^2+|z|^2+\langle |y| \rangle_{\mu}^2),
\end{align*}
for all $t\in [0,T]$, $x,y,z\in\mathbb{R}^n $.
\end{assumption}

 \begin{align}
 U(t, x, \mu)=\ &  V(t, x, \mu) + M(t, \mu) \label{defUU}   \\
   & \ + \langle \mu(dy),  \delta_\mu  V(t, y,\mu; x)- \delta_\mu  V(t, y,\mu; y)      \rangle \notag
 \end{align}

\begin{align}\label{ssp}
\delta_{\mu\mu}\psi(\mu; y,z) =\delta_{\mu\mu} \psi(\mu; z,y)
+\delta_\mu \psi(\mu;y) -\delta_\mu \psi(\mu; z).
\end{align}

\begin{lemma}\label{lemma:xi}
 For \eqref{X1u}--\eqref{Xkuhat},
we have
\begin{align}
& |\xi^a_1 |, |\xi^b|\le \frac{C}{N-1} (1+|X_t^1|^2 +\langle|y|^2\rangle_{\mu_t^{-1}}  ),   \label{xia1b}  \\
& |\xi^a_2 |\le \frac{C}{N-1} ,  \label{xia2}   \\
&|\xi_j^L|\le \frac{C}{|N-1|^2} (1+|X^1_t|^2+|X_t^j|^2+\langle|y|^2\rangle_{\mu_t^{-1}}  ),\label{xiLj}  \\
&|\xi_j^g|\le \frac{C}{|N-1|^2} (1+|X^1_T|^2+|X_T^j|^2+\langle|y|^2\rangle_{\mu_T^{-1}} ).\label{xigj}
\end{align}

\end{lemma}

\begin{proof}
By \eqref{defUU} we calculate $\partial_y \delta_\mu U(t, x, \mu;y)$
and use Assumption \ref{assum:meV} to  obtain
\begin{align}
|\partial_y\delta_\mu U(t, x, \mu;y)|\le C(1+|x|+|y|+\langle|y|\rangle_\mu).
\label{ymuUC}
\end{align}
Next for $\mu^{-j}\coloneqq \frac{1}{N-1}\sum_{k\ne j, k=1}^N \delta_{x^k}$, we have
\begin{align}
&f^*(t, x^j, \mu^{-j})-f^*(t, x^j, \mu^{-1}) \label{ffj1}  \\
&\qquad =
\int_0^1 \int_y\delta_\mu f^*(t, x^j, \mu^{-1} +s(\mu^{-j}-\mu^{-1}) ;y) (\mu^{-j}-\mu^{-1}) (dy)ds. \notag
\end{align}
Denote $\hat \mu_s\coloneqq  \mu^{-1} +s(\mu^{-j}-\mu^{-1})$
and $\Delta^{j,1} \mu\coloneqq \mu^{-j}-\mu^{-1}  $.
  Assumption \ref{assum:flg} implies
\begin{align}
\Big|\int_y \delta_\mu f^*(t, x^j, \hat \mu_s;y)  \Delta^{j,1} \mu (dy)\Big| &=\frac{1}{N-1} |\delta_\mu f^*(t, x^j, \hat \mu_s;x^1)- \delta_\mu f^*(t, x^j, \hat \mu_s;x^j) | \notag \\
&\le \frac{C}{N-1} \Big(1+|x^j|+|x^1|+\frac{1}{N-1}\sum_{k\ge 2} |x^k|\Big).\notag
\end{align}
Subsequently, by \eqref{ymuUC} and \eqref{ffj1}, we have
\begin{align}
|\xi_1^a|&\le \frac{C}{N-1}\Big(1+|X_t^1|^2+\frac{1}{N-1}\sum_{j\ge 2}|X_t^j|^2+ \langle |y|\rangle^2_{\mu_t^{-1}}\Big)\\
&\le \frac{C}{N-1}(1+|X_t^1|^2+ \langle |y|^2\rangle_{\mu_t^{-1}}).\notag
\end{align}

Next, we use the growth conditions of $V$,  as specified by $ {\cal C}_V$, and the semi-symmetry property \eqref{ssp} to establish
\begin{align}
|\partial_{yz}
\delta_{\mu\mu}
U(t, X^1_t,\mu_t^{-1};X_t^k, X^k_t)|\le C,
\end{align}
which implies the bound of $|\xi_2^a|$ in  \eqref{xia2}.
To get the bound of $|\xi^b|$, we write

\begin{align}
f^*(t, x^j, \mu^{-j})=& f^*(t, x^j, \mu^{-1}) +
\langle \Delta^{j,1} \mu(dy)  , \ \delta_\mu f^*(t, x^j, \mu^{-1} ; y)     \rangle \label{fTaylor} \\
& + \int_0^1\int_0^1\int_{y,z} s \delta_{\mu\mu} f^*(t,x^j, \mu^{-1}+ s\tau (\mu^{-j}-\mu^{-1}); y,z) \notag \\
&\qquad\qquad \qquad \cdot  \Delta^{j,1} \mu (dy)
 \Delta^{j,1} \mu  (dz)d\tau  ds.\notag
\end{align}
Denote $\hat\mu_{s,\tau}\coloneqq \mu^{-1}+ s\tau (\mu^{-j}-\mu^{-1}) $. Then
\begin{align*}
&\int_{y,z}\delta_{\mu\mu} f^*(t,x^j,\hat \mu_{s,\tau};y,z)\Delta^{j,1}\mu (dy)\Delta^{j,1}\mu (dz) \\
&=  \frac{1}{N-1} \int_z [\delta_{\mu\mu} f^*(t,x^j,\hat \mu_{s,\tau};x^1,z)- \delta_{\mu\mu} f^*(t,x^j,\hat \mu_{s,\tau};x^j,z) ]\Delta^{j,1}\mu (dz)\\
&=\frac{1}{(N-1)^2}\Big\{ [\delta_{\mu\mu} f^*(t,x^j,\hat \mu_{s,\tau};x^1,x^1)- \delta_{\mu\mu} f^*(t,x^j,\hat \mu_{s,\tau};x^j,x^1) ]\\
&\qquad\qquad\qquad- [\delta_{\mu\mu} f^*(t,x^j,\hat \mu_{s,\tau};x^1,x^j)- \delta_{\mu\mu} f^*(t,x^j,\hat \mu_{s,\tau};x^j,x^j) ]\Big\}.
\end{align*}
Consequently, under Assumption \ref{assum:flg}, we have
\begin{align}
&\Big|\int_{y,z}\delta_{\mu\mu} f^*(t,x^j,\hat \mu_{s,\tau};y,z)\Delta^{j,1}\mu (dy)\Delta^{j,1}\mu (dz)\Big| \label{fsdede}  \\
&\qquad \le \frac{C}{(N-1)^2} (1+|x^1|+|x^j| + \langle|y|\rangle_{\mu^{-1}} ).\notag
\end{align}
We  use Assumption \ref{assum:meV} to obtain
\begin{align}
|\partial_y \delta_\mu \overline U(t, \mu; y)|&\le
|\partial_y V(t, y,\mu)|+
|\langle \mu(dw), \partial_y\delta_\mu V(t, w, \mu; y)  \rangle | \label{ymubU}  \\
&\le C(1+ |y|+\langle |y|\rangle_\mu).\notag
\end{align}
By \eqref{fTaylor}, \eqref{fsdede} and \eqref{ymubU}, we obtain the bound of $|\xi^b|$ in \eqref{xia1b}.
Finally, we follow the method in \eqref{fTaylor} to obtain \eqref{xiLj} and \eqref{xigj}.
\end{proof}

\end{document}
